\documentclass[11pt,a4paper]{article}

\usepackage[utf8]{inputenc}

\usepackage[T1]{fontenc}

\usepackage[polish]{babel}

\usepackage{lmodern}

\usepackage{geometry}

\geometry{margin=27mm}

\usepackage[hidelinks]{hyperref}

\title{Kiedy adaptacyjny interfejs powinien podpowiadać? Ostrożna adaptacja 2D oparta na jakości spojrzenia}

\author{Mateusz Pomianek}

\date{Szkic naukowy, wersja 1.0.0 | 28 listopada 2026}

\begin{document}

\maketitle

\begin{abstract}

Proponowana miara główna to różnica odsetków zbędnych podpowiedzi na wszystkie okazje decyzyjne.

\end{abstract}

\section{Wprowadzenie}

Pomiar spojrzenia na urządzeniach mobilnych dostarcza sygnału o niepewnej trafności (Krafka et al., 2016). Odpowiednie poleganie na automatyzacji zależy od powiązania działania z jakością informacji (Lee \& See, 2004).

\section{Cel i pytania badawcze}

Hipotezy robocze wymagają zatwierdzenia przed rozpoczęciem właściwej weryfikacji.

\paragraph{H1.} Polityka podpowiedzi warunkowana jakością spojrzenia obniży odsetek niepotrzebnych podpowiedzi na wszystkie okazje decyzyjne względem polityki adaptacyjnej bez bramki jakości.

\paragraph{H2.} Odsetek poprawnie zakończonych prób będzie niższy od wyniku trybu statycznego o mniej niż 5 punktów procentowych lub będzie wyższy.

\section{Materiał i metody}

Pierwszorzędowy wskaźnik: różnica odsetków zbędnych podpowiedzi na wszystkie okazje decyzyjne.

Randomizację trzech warunków należy skrzyżować z kolejnością. To, czy podpowiedź była potrzebna, określa stan gry według wcześniej ustalonych kryteriów oceniany bez dostępu do przypisanej polityki. Odsetek pominiętych potrzebnych podpowiedzi pozwoli zinterpretować koszt abstencji.

Przed badaniem z udziałem ludzi należy zatwierdzić wymagania etyczne, zasady zgody i ochronę danych.

\section{Plan analizy}

Zamrożenie estymandu, reguł oceny jakości i obsługi braków nastąpi przed dostępem do danych wynikowych. Raport obejmie oszacowanie kontrastu głównego z przedziałem niepewności i pełne mianowniki. Jeśli przewidziano hipotezy drugorzędowe, przed analizą należy ustalić sposób kontroli wielokrotności.

\section{Wyniki}

W niniejszym pliku nie umieszczono wyników empirycznych, statystyk ani nieprzeprowadzonych analiz. Sekcja wymaga uzupełnienia po uzyskaniu i weryfikacji materiału dowodowego.

\section{Dyskusja i ograniczenia}

Ograniczenia obejmują możliwe błędy pomiaru, dobór próby, alternatywne mechanizmy i ograniczoną możliwość uogólniania. Korelacja nie dowodzi przyczynowości; poprawność techniczna nie jest dowodem znaczenia klinicznego.

\section{Wnioski}

Hipotezy i sposób pomiaru stanowią podstawę oceny wykonalności oraz przygotowania właściwego protokołu.

\section*{Bibliografia}

{\setlength{\parindent}{0pt}
\setlength{\parskip}{5pt}
\hangindent=1.27cm
Krafka, K., Khosla, A., Kellnhofer, P., Kannan, H., Bhandarkar, S., Matusik, W., \& Torralba, A. (2016). Eye tracking for everyone. \textit{Proceedings of the IEEE Conference on Computer Vision and Pattern Recognition}, 2176--2184. \url{https://openaccess.thecvf.com/content_cvpr_2016/html/Krafka_Eye_Tracking_for_CVPR_2016_paper.html}\par
\hangindent=1.27cm
Lee, J. D., \& See, K. A. (2004). Trust in automation: Designing for appropriate reliance. \textit{Human Factors, 46}(1), 50--80. \url{https://doi.org/10.1518/hfes.46.1.50_30392}\par
}

\end{document}
